%%%%%%%%%%%%%%%%%%%%%%%%%%%%%%%%%%%%%%%%%%%%%%%%%%%%%%%%%%%%%%%%%%%%%%%%%%%%%%%
%  MINICURSO: ESTATÍSTICA PARA CIÊNCIA DE DADOS COM PYTHON - EMASC 2026
%  IV Encontro de Matemática de São Cristóvão
%  Liga Acadêmica Matemática & Empresa - UFS
%  Compilar: pdflatex 2x
%  Figuras: figuras/*.png (geradas por make_figs.py)
%%%%%%%%%%%%%%%%%%%%%%%%%%%%%%%%%%%%%%%%%%%%%%%%%%%%%%%%%%%%%%%%%%%%%%%%%%%%%%%
\PassOptionsToPackage{table}{xcolor}
\documentclass[aspectratio=169, 10pt, t]{beamer}
\usepackage[utf8]{inputenc}
\usepackage[T1]{fontenc}
\usepackage[brazil]{babel}
\usepackage{graphicx}
\usepackage{booktabs}
\usepackage{colortbl}
\usepackage{tikz}
\usepackage{hyperref}
\usepackage{csquotes}
\usepackage{amsmath,amssymb}
\usepackage{listings}
\usepackage[most]{tcolorbox}
\usepackage[scaled=0.95]{helvet}
\renewcommand{\familydefault}{\sfdefault}
\usetikzlibrary{calc,shadings,positioning}
\tcbuselibrary{listings}

% ======================== PALETA ========================
\definecolor{emascnight}{RGB}{14, 16, 62}
\definecolor{emascindigo}{RGB}{44, 32, 132}
\definecolor{emascblue}{RGB}{19, 36, 130}
\definecolor{emascsky}{RGB}{3, 145, 191}
\definecolor{emasccyan}{RGB}{52, 224, 200}
\definecolor{emascteal}{RGB}{0, 140, 150}
\definecolor{emascorange}{RGB}{255, 176, 20}
\definecolor{emascamber}{RGB}{226, 120, 0}
\definecolor{emascmist}{RGB}{243, 244, 251}
\definecolor{emascline}{RGB}{208, 211, 232}
\definecolor{textgray}{RGB}{38, 40, 72}
\definecolor{softgray}{RGB}{115, 118, 150}

% ======================== TEMA BASE ========================
\usetheme{default}
\usecolortheme{default}
\setbeamersize{text margin left=1.25cm, text margin right=1.25cm}
\setbeamertemplate{navigation symbols}{}
\setbeamercolor{normal text}{fg=textgray}
\setbeamercolor{structure}{fg=emascindigo}
\setbeamercolor{background canvas}{bg=white}
\setbeamerfont{title}{series=\bfseries, size=\fontsize{27}{31}\selectfont}
\setbeamerfont{subtitle}{series=\mdseries, size=\Large}
\hypersetup{colorlinks=true, urlcolor=emascteal, linkcolor=emascindigo}

% ======================== CONTATO ========================
\newcommand{\ligaemail}{lame.ufs.mat@gmail.com}

% ======================== CÓDIGO PYTHON ========================
\lstdefinestyle{py}{
  language=Python,
  basicstyle=\ttfamily\footnotesize,
  keywordstyle=\bfseries\color{emascindigo},
  stringstyle=\color{emascteal},
  commentstyle=\itshape\color{softgray},
  showstringspaces=false,
  breaklines=true,
  columns=fullflexible,
  keepspaces=true,
  lineskip=1pt,
  literate={á}{{\'a}}1 {é}{{\'e}}1 {í}{{\'i}}1 {ó}{{\'o}}1 {ú}{{\'u}}1
           {â}{{\^a}}1 {ê}{{\^e}}1 {ô}{{\^o}}1 {ã}{{\~a}}1 {õ}{{\~o}}1
           {à}{{\`a}}1 {ç}{{\c c}}1 {Á}{{\'A}}1 {É}{{\'E}}1 {Í}{{\'I}}1
           {Ó}{{\'O}}1 {Ú}{{\'U}}1 {Ã}{{\~A}}1 {Õ}{{\~O}}1 {Ç}{{\c C}}1
}
\newtcblisting{pycode}[1][]{listing only, listing style=py, enhanced, #1,
  colback=emascmist, colframe=emascmist, boxrule=0pt, arc=2.5pt,
  left=8pt, right=8pt, top=5pt, bottom=5pt,
  borderline west={3.5pt}{0pt}{emascsky},
  before skip=4pt, after skip=4pt}
\newcommand{\code}[1]{\texttt{\color{emascindigo}#1}}

% ======================== MACROS DE APOIO ========================
\newcommand{\bigtext}{\spaceskip=0.32em plus 0.12em minus 0.04em\relax\hyphenpenalty=10000\exhyphenpenalty=10000\lefthyphenmin=62\righthyphenmin=62}
\newenvironment{pagetikz}{\tikzpicture[remember picture, overlay, shift={(current page.south west)}]}{\endtikzpicture}

\newcommand{\stagelabel}{}
\newcommand{\setstage}[1]{\renewcommand{\stagelabel}{#1}}
\newcommand{\motive}[1]{\par{\small\itshape\color{softgray}#1}\par\vspace{4pt}}

\newcommand{\darkcanvas}{%
\setbeamertemplate{background canvas}{%
\begin{pagetikz}
\shade[shading=axis, shading angle=35, top color=emascnight, bottom color=emascindigo!85!emascnight]
(0,0) rectangle (16,9);
\draw[emasccyan, opacity=0.10, line width=0.7pt] (14.2,7.6) circle (5.2);
\draw[emasccyan, opacity=0.07, line width=0.7pt] (14.2,7.6) circle (6.6);
\draw[emascorange, opacity=0.10, line width=0.7pt] (14.2,7.6) circle (3.8);
\end{pagetikz}}%
}

\newcommand{\lammotif}[5]{%
\begin{scope}[shift={(#1,#2)}, scale=#3, opacity=#5]
\draw[#4, line width=1.1pt, line join=round]
(0,0.1)--(0.8,1.5)--(1.6,0.3)--(2.5,1.7)--(3.3,0.6)--(4.2,1.3);
\foreach \x/\y in {0/0.1,0.8/1.5,1.6/0.3,2.5/1.7,3.3/0.6,4.2/1.3}
\draw[#4, line width=1pt, fill=emascnight] (\x,\y) circle (0.1);
\end{scope}}

\newcommand{\bigsigma}[4]{%
\begin{scope}[shift={(#1,#2)}, scale=#3, opacity=#4]
\draw[white, line width=7pt, line join=round, line cap=round]
(1,1.9)--(0,1.9)--(0.95,0.95)--(0,0)--(1,0);
\end{scope}}

% ======================== BLOCOS ========================
\newcommand{\emascbox}[3]{%
\par\vspace{0.6ex}%
\begin{tcolorbox}[enhanced, breakable=false, colback=emascmist, colframe=emascmist, boxrule=0pt,
arc=2.5pt, left=11pt, right=8pt, top=5pt, bottom=5pt, boxsep=1pt,
borderline west={3.5pt}{0pt}{#1}, shadow={0.6mm}{-0.6mm}{0mm}{black!6}]
\ifx\relax#2\relax\else{\color{#1}\bfseries\small #2}\par\vspace{2pt}\fi
}
\setbeamertemplate{block begin}{\emascbox{emascindigo}{\insertblocktitle}{}}
\setbeamertemplate{block end}{\end{tcolorbox}\vspace{0.3ex}}
\setbeamertemplate{block alerted begin}{\emascbox{emascorange}{\insertblocktitle}{}}
\setbeamertemplate{block alerted end}{\end{tcolorbox}\vspace{0.3ex}}
\setbeamertemplate{block example begin}{\emascbox{emasccyan!80!emascteal}{\insertblocktitle}{}}
\setbeamertemplate{block example end}{\end{tcolorbox}\vspace{0.3ex}}

% ======================== LISTAS ========================
\setbeamertemplate{itemize item}{\tikz[baseline=-0.5ex]\fill[emascorange] (0,0) circle (0.085cm);}
\setbeamertemplate{itemize subitem}{\tikz[baseline=-0.5ex]\draw[emascteal, line width=0.9pt] (0,0) circle (0.07cm);}
\setbeamertemplate{itemize subsubitem}{\tikz[baseline=-0.5ex]\fill[emascline] (0,0) circle (0.05cm);}
\setbeamertemplate{enumerate item}{%
\tikz[baseline=(n.base)]
\node[circle, fill=emascindigo, text=white, inner sep=1.4pt, minimum size=0.42cm,
font=\scriptsize\bfseries] (n) {\insertenumlabel};}
\setbeamertemplate{enumerate subitem}{\insertenumlabel.\insertsubenumlabel}
\setlength{\leftmargini}{1.6em}
\setbeamertemplate{caption label separator}{: }
\setbeamerfont{caption}{size=\small}
\setbeamercolor{caption name}{fg=emascindigo}

% ======================== CABEÇALHO ========================
\setbeamertemplate{frametitle}{%
\vbox to 1.5cm{%
\begin{tikzpicture}[remember picture, overlay, shift={(current page.north west)}]
\node[anchor=west, text=emascindigo, font=\fontsize{19}{22}\bfseries\selectfont, text width=12cm]
at (1.25,-0.82) {\insertframetitle};
\fill[emascorange, rounded corners=0.5pt] (1.25,-1.38) rectangle (2.55,-1.43);
\shade[left color=emasccyan, right color=emasccyan!0] (2.6,-1.38) rectangle (10.5,-1.43);
\node[anchor=east] at (14.85,-0.82) {\includegraphics[height=1.05cm]{figuras/logo_liga.png}};
\end{tikzpicture}\vfil}%
}

% ======================== RODAPÉ ========================
\setbeamertemplate{footline}{%
\vbox to 0.8cm{\begin{tikzpicture}[remember picture, overlay, shift={(current page.south west)}]
\draw[emascline, line width=0.4pt] (1.25,0.62)--(14.75,0.62);
\node[anchor=west, font=\fontsize{7}{8}\selectfont, text=softgray] at (1.25,0.33)
{\textbf{EMASC 2026}\,\;\textcolor{emascorange}{\textbullet}\;\,Liga Acadêmica Matemática \& Empresa\,\;\textcolor{emascorange}{\textbullet}\;\,UFS};
\node[anchor=east, font=\fontsize{7}{8}\selectfont, text=softgray] at (14.75,0.33)
{\textbf{\textcolor{emascindigo}{\insertframenumber}}\,/\,\inserttotalframenumber};
\shade[shading=axis, shading angle=90, top color=emascindigo, bottom color=emascsky] (0,0) rectangle (0.2,9);
\fill[emascorange] (0,8.15) rectangle (0.2,9);
\end{tikzpicture}\vfil}%
}

% ======================== SUMÁRIO ========================
\setbeamertemplate{section in toc}{%
\leavevmode\vskip 1.1ex
\tikz[baseline=(n.base)]
\node[circle, fill=emascindigo, text=white, inner sep=2pt, minimum size=0.55cm,
font=\small\bfseries] (n) {\inserttocsectionnumber};%
\hspace{0.7em}{\large\color{textgray}\inserttocsection}\par}
\setbeamertemplate{subsection in toc}{}
\setbeamercolor{section in toc}{fg=textgray}

% ======================== SLIDE DE SEÇÃO ========================
\AtBeginSection[]{%
{\darkcanvas
\begin{frame}[plain, noframenumbering]
\begin{pagetikz}
\lammotif{9.2}{1.2}{1.35}{emasccyan}{0.55}
\bigsigma{12.4}{4.6}{2.2}{0.05}
\node[anchor=south west, text=emasccyan, font=\fontsize{9}{11}\bfseries\selectfont]
at (1.45,7.4) {\stagelabel};
\node[anchor=west, text=emascorange, font=\fontsize{60}{60}\bfseries\selectfont]
at (1.4,5.6) {\ifnum\value{section}<10 0\fi\arabic{section}};
\fill[emasccyan] (1.45,4.55) rectangle (3.2,4.6);
\node[anchor=north west, text=white, font=\fontsize{26}{30}\bfseries\selectfont, text width=12.8cm, align=left]
at (1.4,4.3) {\bigtext\hypersetup{linkcolor=white}\insertsection\par};
\node[anchor=south west, text=emasccyan, font=\fontsize{8}{10}\bfseries\selectfont]
at (1.4,0.6) {EMASC 2026\;\textcolor{emascorange}{\textbullet}\;LAME};
\end{pagetikz}
\end{frame}}}

% ======================== METADADOS ========================
\title[Estatística com Python]{Estatística para Ciência de Dados com Python}
\institute[LAME - UFS]{%
Liga Acadêmica Matemática \& Empresa \\
Universidade Federal de Sergipe -- UFS
}
\date[2026]{São Cristóvão -- SE, 2026}

% ======================== MACROS DO MINICURSO ========================
\newcommand{\figbox}[2]{\parbox[c][#1][c]{\linewidth}{\centering\includegraphics[width=\linewidth,height=#1,keepaspectratio]{figuras/#2}}}

\newcommand{\tip}[2]{%
\begin{tcolorbox}[enhanced, colback=emascmist, colframe=emascmist, boxrule=0pt,
arc=2.5pt, left=11pt, right=9pt, top=5pt, bottom=5pt,
borderline west={3.5pt}{0pt}{#1}, fontupper=\small,
before skip=4pt, after skip=2pt]#2\end{tcolorbox}}

% \card{cor}{grupo de altura (único por slide)}{título}{corpo}
\newcommand{\card}[4]{%
\begin{tcolorbox}[enhanced, colback=emascmist, colframe=#1, boxrule=0pt, toprule=3.5pt,
arc=3pt, left=7pt, right=7pt, top=6pt, bottom=6pt,
equal height group=#2, valign=top, before skip=2pt, after skip=2pt]
{\bfseries\color{#1!80!black}#3}\par\vspace{3pt}
{\raggedright #4\par}
\end{tcolorbox}}

\newcommand{\pnode}[1][emascsky]{
\node[draw=#1, fill=#1!15, rounded corners=3pt, inner sep=4pt, font=\small\bfseries]}

\newcommand{\tp}[1]{\par\vspace{1pt}\textcolor{emascorange}{\textbullet}\,#1}
\newcommand{\passo}[2]{\textcolor{emascindigo}{\textbf{#1.}}\ #2\par\vspace{2pt}}
\newcommand{\dmath}[1]{{\small\[#1\]}}
\newcommand{\dmathf}[1]{{\footnotesize\[#1\]}}

\begin{document}

% -------------------- TÍTULO --------------------
{\darkcanvas
\begin{frame}[plain, noframenumbering]
\begin{pagetikz}
\bigsigma{11.3}{3.6}{3.1}{0.06}
\lammotif{10.6}{6.9}{1.0}{emasccyan}{0.7}
\lammotif{10.9}{7.25}{1.0}{emascorange}{0.35}
\node[anchor=west, fill=emascorange, text=emascnight, rounded corners=7pt, inner xsep=9pt, inner ysep=3.5pt, font=\fontsize{8}{9}\bfseries\selectfont] at (1.25,8.2) {EMASC\,2026};
\node[anchor=west, text=emasccyan, font=\fontsize{8}{9}\bfseries\selectfont] at (3.7,8.2) {LIGA ACADÊMICA MATEMÁTICA \& EMPRESA};
\node[anchor=north west, text=white, text width=9.0cm, align=left] at (1.25,7.7) {\bigtext\fontsize{22}{26}\selectfont\bfseries\inserttitle\par};
\fill[emascorange] (1.3,4.7) rectangle (2.9,4.76);
\fill[emasccyan] (3.0,4.7) rectangle (3.6,4.76);
\node[anchor=north west, text=emasccyan, font=\fontsize{14}{17}\bfseries\selectfont] at (1.25,4.5) {\insertsubtitle};
\node[anchor=north west, text=white, font=\fontsize{13}{16}\bfseries\selectfont] at (1.25,3.65) {\insertauthor};
\node[anchor=north west, text=white!78!emascindigo, font=\fontsize{8.5}{11}\selectfont, align=left] at (1.25,3.15) {\insertinstitute};
\node[anchor=north west, text=emascorange, font=\fontsize{8}{10}\bfseries\selectfont] at (1.25,2.2) {\insertdate};
\fill[white] (0,0) rectangle (16,1.65);
\fill[emascorange] (0,1.65) rectangle (5.3,1.69);
\shade[left color=emasccyan, right color=emascsky] (5.3,1.65) rectangle (16,1.69);
\node at (4.6,0.83)  {\includegraphics[height=1.25cm]{figuras/logo_ufs.png}};
\node at (8.0,0.83)  {\includegraphics[height=1.3cm]{figuras/logo_emasc.png}};
\node at (11.4,0.83) {\includegraphics[height=1.25cm]{figuras/logo_liga.png}};
\draw[emascline, line width=0.5pt] (6.3,0.35)--(6.3,1.3);
\draw[emascline, line width=0.5pt] (9.7,0.35)--(9.7,1.3);
\end{pagetikz}
\end{frame}}



% ============================================================
\setstage{INTRODUÇÃO}
\section{Introdução à estatística}

\begin{frame}{O que é estatística?}
\motive{Coletar, organizar, resumir e interpretar dados, quantificando a incerteza das conclusões.}
\begin{columns}[T]
\begin{column}{.48\textwidth}\card{emascsky}{i1}{Descritiva}{Resume e descreve os dados que temos: médias, tabelas e gráficos.\par\vspace{6pt}{\small\bfseries\color{emascamber} É o foco deste minicurso.}}\end{column}
\begin{column}{.48\textwidth}\card{softgray}{i1}{Inferencial}{Usa uma \textbf{amostra} para concluir sobre um grupo maior, com uma medida de confiança.\par\vspace{6pt}{\small\color{softgray}Intervalos de confiança e testes de hipótese.}}\end{column}
\end{columns}
\end{frame}

\begin{frame}[fragile]{Ferramentas do minicurso}
\motive{Quatro bibliotecas. A SciPy é a referência dos cálculos; pandas e NumPy têm equivalentes.}
\begin{tcolorbox}[enhanced, colback=emascorange!12, colframe=emascorange,
boxrule=1.2pt, toprule=3.5pt, arc=3pt, left=9pt, right=9pt, top=5pt, bottom=5pt,
lower separated=true, colbacklower=emascorange!12, segmentation style={emascorange!55, dashed, line width=0.6pt}, before skip=2pt, after skip=2pt]
{\bfseries\large\color{emascamber}SciPy}\hfill{\small\color{emascamber}módulo \code{scipy.stats}}\par\vspace{3pt}
{\footnotesize
\begin{minipage}[t]{.52\linewidth}
\code{tmean}, \code{tvar}, \code{tstd}: média, variância, desvio padrão\par
\code{scoreatpercentile}, \code{iqr}: mediana, quartis e IQR\par
\code{mode}, \code{skew}, \code{zscore}: moda, assimetria, padronização\end{minipage}\hfill
\begin{minipage}[t]{.44\linewidth}
\code{binned\_statistic}: contagem por faixa\par
\code{pearsonr}: correlação\par
\code{norm}, \code{binom}, \code{poisson}: distribuições\end{minipage}}
\tcblower
{\bfseries\color{emascsky!80!black}Apoio}\par\vspace{3pt}
{\footnotesize
\begin{minipage}[t]{.31\linewidth}\code{pandas}: ler a tabela de dados\end{minipage}\hfill
\begin{minipage}[t]{.31\linewidth}\code{NumPy}: vetores e máscaras\end{minipage}\hfill
\begin{minipage}[t]{.31\linewidth}\code{Matplotlib}: gráficos\end{minipage}}
\end{tcolorbox}
\begin{pycode}
import pandas as pd, numpy as np, matplotlib.pyplot as plt
from scipy import stats
\end{pycode}
\end{frame}

\begin{frame}[fragile]{A base de dados: a população}
\motive{Um CSV com $N=16.000$ alunos e 10 variáveis. Ele é a nossa \textbf{população}.}
\renewcommand{\arraystretch}{0.92}
\centering\footnotesize
\begin{tabular}{lll}
\toprule
\textbf{Variável} & \textbf{Descrição} & \textbf{Tipo} \\
\midrule
\code{ID\_Aluno} & identificador do aluno & identificador \\
\code{Idade} & idade em anos (13 a 25) & quantitativa discreta \\
\code{Genero} & Masculino, Feminino ou Não-binário & qualitativa nominal \\
\code{Escolaridade} & Ensino Médio, Faculdade ou Universidade & qualitativa ordinal \\
\code{Horas\_Redes} & horas por dia em redes sociais & quantitativa contínua \\
\code{Horas\_IA} & horas por dia em ferramentas de IA & quantitativa contínua \\
\code{Horas\_Sono} & horas de sono por dia & quantitativa contínua \\
\code{Horas\_Atividade} & horas de atividade física por dia & quantitativa contínua \\
\code{Saude\_Mental} & pontuação de saúde mental & quantitativa contínua \\
\code{Saude\_Fisica} & pontuação de saúde física & quantitativa contínua \\
\bottomrule
\end{tabular}
\begin{pycode}
df = pd.read_csv("alunos_redes_ia.csv")   # df.shape: (16000, 10)
\end{pycode}
\end{frame}

\begin{frame}[fragile]{População e amostra}
\motive{Quase nunca observamos a população inteira: usamos uma amostra para estimá-la.}
\begin{columns}[c]
\begin{column}{.46\textwidth}
\centering
\begin{tikzpicture}[scale=0.92]
\fill[emascmist] (0,0) ellipse (2.5 and 1.2); \draw[emascsky, very thick] (0,0) ellipse (2.5 and 1.2);
\foreach \px/\py in {0.6/0.91,1.32/-0.63,-0.96/0.86,1.43/-0.07,-0.95/-0.51,-1.18/-0.13,0.02/0.12,-1.37/-0.78,0.54/-1.05,-2.23/0.03,-0.16/0.96,0.62/0.03,-0.02/-0.58,0.92/-0.69,1.58/-0.79,-1.12/0.87,0.05/0.8,0.67/0.56,-1.96/0.09,0.04/0.85,-0.67/0.23,-2.12/-0.26,-0.85/-0.8,1.52/-0.28,2.3/0.21,0.5/0.32,0.85/-0.8,-0.29/-0.6,-0.47/-0.93,0.82/-0.46,1.8/0.37,0.33/-0.82,0.25/-0.73,1.84/0.33,0.33/-0.28,-0.43/-0.6,-0.15/0.11,-0.85/0.58,-2.28/-0.29,2.24/0.36,-0.34/0.05,1.79/-0.36,0.43/0.42,-0.69/0.04,1.27/0.94,1.49/-0.84}{\fill[emascsky!55] (\px,\py) circle (0.045);}
\foreach \px/\py in {-1.3/0.55,-0.5/-0.2,0.2/0.55,0.75/-0.35,1.35/0.2,-0.9/-0.65,0.1/-0.7,1.0/0.7,-1.7/-0.1,-0.2/0.1,0.5/0.05,1.7/-0.4}{\fill[emascorange] (\px,\py) circle (0.095);}
\node[font=\small\bfseries, text=emascsky] at (0,1.55) {população: $N=16.000$};
\node[font=\small\bfseries, text=emascamber] at (0,-1.6) {amostra: $n=12$};
\end{tikzpicture}
\end{column}
\begin{column}{.50\textwidth}
\small\centering
\begin{tabular}{lcc}
\toprule
 & \textbf{população} & \textbf{amostra} \\
\midrule
tamanho & $N=16.000$ & $n=12$ \\
média de \code{Horas\_Redes} & $\mu=4{,}54$ & $\bar{x}=4{,}94$ \\
desvio padrão & $\sigma=2{,}40$ & $s=2{,}27$ \\
\midrule
\textbf{nome} & parâmetro & estatística \\
\bottomrule
\end{tabular}
\end{column}
\end{columns}
\vspace{2pt}
\begin{pycode}
amostra = df.sample(12, random_state=743)   # 12 alunos sorteados
\end{pycode}
\end{frame}

\begin{frame}{A amostra: 12 alunos sorteados}
\motive{Cada linha é um aluno sorteado ao acaso (semente fixa) do CSV. Os cálculos à mão usarão estes 12.}
\renewcommand{\arraystretch}{0.88}
\centering\footnotesize
\begin{tabular}{lcllcc}
\toprule
\code{ID\_Aluno} & \code{Idade} & \code{Genero} & \code{Escolaridade} & \code{Horas\_Redes} & \code{Saude\_Mental} \\
\midrule
ALU\_05711 & 16 & Masculino & Ensino Médio & 5,00 & 77,85 \\
ALU\_11704 & 22 & Feminino & Universidade & 4,91 & 76,81 \\
ALU\_11034 & 20 & Feminino & Faculdade & 7,10 & 76,53 \\
ALU\_11755 & 21 & Não-binário & Faculdade & 9,93 & 67,72 \\
ALU\_01744 & 13 & Feminino & Ensino Médio & 4,33 & 69,51 \\
ALU\_07422 & 16 & Feminino & Ensino Médio & 4,71 & 80,37 \\
ALU\_10529 & 18 & Masculino & Faculdade & 7,39 & 81,97 \\
ALU\_02039 & 24 & Feminino & Universidade & 3,61 & 72,18 \\
ALU\_01124 & 13 & Masculino & Ensino Médio & 4,12 & 80,17 \\
ALU\_01781 & 16 & Masculino & Ensino Médio & 4,17 & 82,05 \\
ALU\_02360 & 20 & Masculino & Faculdade & 1,92 & 84,26 \\
ALU\_11449 & 21 & Feminino & Faculdade & 2,04 & 81,11 \\
\bottomrule
\end{tabular}
\end{frame}

\begin{frame}{O que é uma variável?}
\motive{Variável: uma característica medida em cada indivíduo, que muda de um para outro.}
\begin{columns}[c]
\begin{column}{.62\textwidth}
\centering
\begin{tikzpicture}[scale=0.76, every node/.style={font=\small}]
\fill[emascorange!28] (5.6,0.3) rectangle (7.5,-3.3);
\fill[emascsky!22] (0,-1.5) rectangle (7.5,-2.1);
\fill[emascorange!80] (5.6,-1.5) rectangle (7.5,-2.1);
\draw[emascamber, very thick] (5.6,-1.5) rectangle (7.5,-2.1);
\draw[softgray] (0,0.3)--(7.5,0.3); \draw[softgray] (0,-0.3)--(7.5,-0.3); \draw[softgray] (0,-3.3)--(7.5,-3.3);
\node[font=\footnotesize\bfseries] at (1.15,0) {ID\_Aluno};
\node[font=\footnotesize\bfseries] at (2.85,0) {Idade};
\node[font=\footnotesize\bfseries] at (4.4,0) {Genero};
\node[font=\footnotesize\bfseries] at (6.55,0) {Horas\_Redes};
\node[font=\footnotesize] at (1.15,-0.6) {ALU\_05711}; \node[font=\footnotesize] at (2.85,-0.6) {16}; \node[font=\footnotesize] at (4.4,-0.6) {Masculino}; \node[font=\footnotesize] at (6.55,-0.6) {5,00};
\node[font=\footnotesize] at (1.15,-1.2) {ALU\_11704}; \node[font=\footnotesize] at (2.85,-1.2) {22}; \node[font=\footnotesize] at (4.4,-1.2) {Feminino}; \node[font=\footnotesize] at (6.55,-1.2) {4,91};
\node[font=\footnotesize] at (1.15,-1.8) {ALU\_11034}; \node[font=\footnotesize] at (2.85,-1.8) {20}; \node[font=\footnotesize] at (4.4,-1.8) {Feminino}; \node[font=\footnotesize] at (6.55,-1.8) {7,10};
\node[font=\footnotesize] at (1.15,-2.4) {ALU\_11755}; \node[font=\footnotesize] at (2.85,-2.4) {21}; \node[font=\footnotesize] at (4.4,-2.4) {Não-binário}; \node[font=\footnotesize] at (6.55,-2.4) {9,93};
\foreach \xx in {1.15,2.85,4.4,6.55} \node at (\xx,-3.0) {$\vdots$};
\node[text=emascamber, font=\small\bfseries] at (6.55,1.15) {variável (coluna)};
\draw[->, emascamber, thick] (6.55,0.85)--(6.55,0.4);
\node[text=emascsky, font=\small\bfseries, anchor=west, align=left] at (8.0,-0.9) {observação\\(linha = 1 aluno)};
\draw[->, emascsky, thick] (7.95,-1.2)--(7.58,-1.7);
\node[text=emascamber, font=\small\bfseries, anchor=west, align=left] at (8.0,-3.0) {valor\\(célula)};
\draw[->, emascamber, thick] (7.95,-2.75)--(7.55,-2.0);
\end{tikzpicture}
\end{column}
\begin{column}{.34\textwidth}
\small
\passo{1}{\textbf{Variável}: cada coluna, como \code{Horas\_Redes}.}
\passo{2}{\textbf{Observação}: cada linha, um aluno.}
\passo{3}{\textbf{Valor}: o dado na célula, como 7,10 h.}
\end{column}
\end{columns}
\tip{emascorange}{Uma coluna é variável \textbf{porque o valor muda} de aluno para aluno. \code{ID\_Aluno} só identifica cada linha.}
\end{frame}

\begin{frame}{Dados qualitativos e quantitativos}
\motive{O tipo da variável decide quais contas e quais gráficos fazem sentido.}
\tip{emascsky}{\textbf{Qualitativas}: categorias (nominal ou ordinal). \textbf{Quantitativas}: números (discreta ou contínua).}
\begin{columns}[T]
\begin{column}{.235\textwidth}\card{emascorange}{v1}{Nominal}{\centering\tikz{\foreach \i/\c in {0/emascorange,1/emascsky,2/emasccyan}\fill[\c] (\i*0.45,0) circle(0.17);}\par\vspace{5pt}\raggedright\scriptsize\textbf{Genero}\par categorias sem ordem}\end{column}
\begin{column}{.235\textwidth}\card{emascsky}{v1}{Ordinal}{\centering\tikz{\fill[emascsky!35](0,0)rectangle(.28,.15);\fill[emascsky!65](.35,0)rectangle(.63,.3);\fill[emascsky](.7,0)rectangle(.98,.45);}\par\vspace{5pt}\raggedright\scriptsize\textbf{Escolaridade}\par categorias com ordem}\end{column}
\begin{column}{.235\textwidth}\card{emascteal}{v1}{Discreta}{\centering\tikz{\draw[softgray](0,0)--(1.3,0);\foreach \i in {0,...,4}\fill[emascteal](\i*0.32,0)circle(0.07);}\par\vspace{5pt}\raggedright\scriptsize\textbf{Idade}\par contagens: valores inteiros}\end{column}
\begin{column}{.235\textwidth}\card{emascindigo}{v1}{Contínua}{\centering\tikz{\shade[left color=emasccyan,right color=emascindigo](0,0)rectangle(1.3,.2);}\par\vspace{5pt}\raggedright\scriptsize\textbf{Horas\_Redes}\par medidas em escala real}\end{column}
\end{columns}
\vspace{0.1cm}
\end{frame}

\begin{frame}[fragile]{Tipos de variável em Python}
\motive{No pandas, o tipo de cada coluna aparece no \code{dtype}. A ordem de uma ordinal precisa ser declarada.}
\begin{pycode}
df.dtypes                          # tipo de cada coluna
ordem = ["Ensino Médio", "Faculdade", "Universidade"]
esc = pd.Categorical(df["Escolaridade"], categories=ordem, ordered=True)
esc.min()                          # 'Ensino Médio': agora há uma ordem
\end{pycode}
\tip{emascorange}{Um tipo errado (por exemplo, número lido como texto) é um dos erros mais comuns na prática.}
\end{frame}

\begin{frame}[fragile]{Por que explorar antes de modelar?}
\motive{Um modelo treina com qualquer dado e devolve um número, sem avisar se algo estava errado.}
\begin{columns}[T]
\begin{column}{.235\textwidth}\card{emascsky}{x1}{1. Qualidade}{\small Nulos, duplicatas, tipos errados.}\end{column}
\begin{column}{.235\textwidth}\card{emascteal}{x1}{2. Formato}{\small Concentração, simetria e extremos.}\end{column}
\begin{column}{.235\textwidth}\card{emascindigo}{x1}{3. Método}{\small Outliers e assimetria mudam a escolha.}\end{column}
\begin{column}{.235\textwidth}\card{emascorange}{x1}{4. Hipóteses}{\small Padrões orientam o que testar.}\end{column}
\end{columns}
\vspace{0.1cm}
\begin{pycode}
df.shape                 # (16000, 10)
df.head()                # primeiras linhas
df.info()                # tipos e nulos por coluna
df.isna().sum().sum()    # 0 nulos
\end{pycode}
\end{frame}

% ============================================================
\setstage{ESTATÍSTICA DESCRITIVA}
\section{Estatística descritiva}

\begin{frame}{Média aritmética}
\motive{O ponto de equilíbrio dos dados. Cada bolinha é um aluno da amostra (horas por dia em redes sociais).}
\figbox{3.9cm}{gangorra_a}
\begin{columns}[c]
\begin{column}{.56\textwidth}
\dmath{\bar{x}=\frac{1}{n}\sum_{i=1}^{n}x_i=\frac{1{,}92+2{,}04+\cdots+9{,}93}{12}=\frac{59{,}23}{12}\approx4{,}94\ \text{h}}
\end{column}
% \begin{column}{.40\textwidth}\figbox{1.3cm}{puxoes_a}\end{column}
\end{columns}
\end{frame}

\begin{frame}{Mediana}
\motive{O valor central dos dados ordenados.}
\figbox{3.0cm}{mediana}\par\vspace{2pt}
\begin{columns}[T]
\begin{column}{.52\textwidth}
\passo{1}{Ordenar os dados (a figura já está ordenada).}
\passo{2}{$n=12$ é par: posições $6$ e $7$.}
\passo{3}{Média dos dois centrais: $\frac{4{,}33+4{,}71}{2}=4{,}52$.}
\end{column}
\begin{column}{.44\textwidth}
\small
Se $n$ é \textbf{ímpar}: mediana $=x_{((n+1)/2)}$. A mediana é o \textbf{percentil 50}: metade dos dados fica abaixo dela.
\end{column}
\end{columns}
\end{frame}

\begin{frame}{Moda}
\motive{O valor mais frequente. Serve bem para variáveis categóricas e discretas.}
\figbox{3.4cm}{moda}\par\vspace{2pt}
\small
\passo{1}{\textbf{Genero}: Feminino aparece 6 vezes, Masculino 5 e Não-binário 1. Moda: \textbf{Feminino}.}
\passo{2}{\textbf{Idade}: 16 anos aparece 3 vezes. Moda: \textbf{16}.}
\passo{3}{\textbf{Horas\_Redes}: nenhum valor se repete, então não há moda: em dados contínuos ela ajuda pouco.}
\end{frame}

\begin{frame}{Comparando média e mediana}
\only<1>{\motive{Apoiando na mediana (4,52), o 9,93 pesa mais e a prancha desce à direita.}}
\only<2>{\motive{A média (4,94) é o ponto que equilibra: o 9,93 puxa o apoio para a direita.}}
\only<3>{\motive{Sem o aluno de 9,93 h: média e mediana ficam próximas.}}
\only<1>{\figbox{4.1cm}{gangorra_med}}\only<2>{\figbox{4.1cm}{gangorra_a}}\only<3>{\figbox{4.1cm}{gangorra_b}}\par\vspace{2pt}
\only<1>{\tip{emascsky}{A mediana divide a \textbf{quantidade} de alunos (6 de cada lado), mas ignora \textbf{quão longe} eles estão.}}
\only<2>{\tip{emascorange}{Média $>$ mediana (4,94 $>$ 4,52) indica \textbf{assimetria positiva}: uma cauda longa à direita. O coeficiente de assimetria é $0{,}77$.}}
\only<3>{\tip{emascsky}{Retirando o 9,93 ($n=11$): média $\approx4{,}48$ e mediana $=4{,}33$.}}
\end{frame}

\begin{frame}[fragile]{Em Python: posição e assimetria}
\begin{pycode}
x = amostra["Horas_Redes"].to_numpy()

stats.tmean(x)                      # 4.936  (média)
stats.scoreatpercentile(x, 50)      # 4.52   (mediana = percentil 50)
stats.mode(amostra["Idade"]).mode   # 16     (moda da idade)
amostra["Genero"].mode()[0]         # 'Feminino' (moda do gênero)
stats.skew(x)                       # 0.767  (assimetria)
\end{pycode}
\tip{emascsky}{O \code{scipy.stats} não tem \code{median}: usa-se o percentil 50. \code{stats.mode} só aceita números; para texto, use o \code{mode()} do pandas.}
\end{frame}

\begin{frame}{Quartis}
\only<1>{\motive{Os quartis são percentis: $Q_1$ (25\%), mediana (50\%) e $Q_3$ (75\%). Ficam entre dois valores vizinhos.}}
\only<2>{\motive{Interpolar: partir do valor anterior e avançar a fração da distância até o seguinte.}}
\only<1>{\figbox{3.4cm}{quartis1}}\only<2>{\figbox{3.4cm}{quartis2}}\par\vspace{2pt}
\only<1>{\dmath{Q_1=\operatorname{mediana}\bigl(x_{(1)},\dots,x_{(n/2)}\bigr),\qquad Q_3=\operatorname{mediana}\bigl(x_{(n/2+1)},\dots,x_{(n)}\bigr)}}
\only<2>{\tip{emascorange}{$Q_1$: $h=1+11\cdot0{,}25=3{,}75$. $Q_3$: $h=1+11\cdot0{,}75=9{,}25$. Este é o método padrão do SciPy, do NumPy e do \code{plt.boxplot}.}}
\end{frame}

\begin{frame}[fragile]{Em Python: mínimo, máximo e quartis}
\begin{pycode}
x = amostra["Horas_Redes"].to_numpy()

stats.tmin(x), stats.tmax(x)               # 1.92, 9.93
stats.scoreatpercentile(x, [25, 50, 75])   # [3.9925, 4.52, 5.525]
stats.iqr(x)                               # 1.5325 (Q3 - Q1)
\end{pycode}
\tip{emascsky}{\code{scoreatpercentile} recebe o percentil em porcentagem (0 a 100). \code{stats.iqr} já devolve $Q_3-Q_1$.}
\end{frame}

\begin{frame}{Variância e desvio padrão}
\motive{Quanto, em média, cada valor se afasta da média?}
\dmath{s^2=\frac{1}{n-1}\sum_{i=1}^{n}(x_i-\bar{x})^2,\qquad s=\sqrt{s^2}}
\begin{columns}[T]
\begin{column}{.48\textwidth}
\passo{1}{Calcular a média $\bar{x}$.}
\passo{2}{Calcular os desvios $x_i-\bar{x}$.}
\passo{3}{Elevar ao quadrado e somar.}
\end{column}
\begin{column}{.48\textwidth}
\passo{4}{Dividir por $n-1$: obtém-se $s^2$.}
\passo{5}{Extrair a raiz: $s$, na unidade original.}
\end{column}
\end{columns}
\tip{emascorange}{Por que $n-1$? Os desvios em relação a $\bar{x}$ somam 0, então só $n-1$ deles são livres. Dividir por $n$ subestimaria a variância da população.}
\end{frame}

\begin{frame}{Desvios e quadrados}
\only<1>{\figbox{4.4cm}{desvios}\par\vspace{2pt}\tip{emascsky}{Com $\bar{x}=4{,}94$: cada $d_i=x_i-\bar{x}$ é a distância com sinal até a média, e $\sum d_i=0$. O desvio de $+4{,}99$ domina.}}
\only<2>{\figbox{4.4cm}{quadrados}\par\vspace{2pt}\dmath{\sum(x_i-\bar{x})^2=56{,}56,\quad s^2=\frac{56{,}56}{11}\approx5{,}14,\quad s=\sqrt{5{,}14}\approx2{,}27\ \text{h}}}
\end{frame}

\begin{frame}[fragile]{Em Python: variância e desvio}
\begin{pycode}
x = amostra["Horas_Redes"].to_numpy()

d = x - stats.tmean(x)             # desvios
(d**2).sum()                       # 56.56
stats.tvar(x)                      # 5.142  (variância amostral, n-1)
stats.tstd(x)                      # 2.268  (desvio padrão)
stats.variation(x, ddof=1)         # 0.459  (coeficiente de variação)
\end{pycode}
\tip{emascorange}{\code{tvar} e \code{tstd} já dividem por $n-1$. Em \code{variation} é preciso escrever \code{ddof=1}, pois o padrão é \code{ddof=0}.}
\end{frame}

\begin{frame}{Outliers: a regra do IQR}
\motive{Um outlier é um valor que foge do padrão dos demais.}
\figbox{3.7cm}{limites}\par\vspace{2pt}
\begin{columns}[T]
\begin{column}{.56\textwidth}
\dmathf{L_{\inf}=Q_1-1{,}5\,\mathrm{IQR}=3{,}99-2{,}30=1{,}69}
\dmathf{L_{\sup}=Q_3+1{,}5\,\mathrm{IQR}=5{,}53+2{,}30=7{,}82}
\end{column}
\begin{column}{.40\textwidth}
\small Como $9{,}93>7{,}82$, esse aluno é discrepante.
\end{column}
\end{columns}
\end{frame}

\begin{frame}{Padronização: o escore $z$}
\motive{Quantos desvios padrão um valor está acima ou abaixo da média?}
\figbox{3.3cm}{zscore}\par\vspace{0pt}
\dmath{z_i=\frac{x_i-\bar{x}}{s},\qquad z=\frac{9{,}93-4{,}94}{2{,}27}\approx2{,}20}
\tip{emascsky}{Regra comum: $|z|>3$ é discrepante. Mas, com $n=12$, o maior $z$ possível é $\frac{n-1}{\sqrt{n}}\approx3{,}18$: a regra quase nunca dispara. Por isso a regra do IQR é mais útil em amostras pequenas.}
\end{frame}

\begin{frame}{O efeito de um único valor extremo}
\motive{Retirando o 9,93, observe o que muda e o que permanece.}
\figbox{3.9cm}{comparacao}\par\vspace{2pt}
\tip{emascorange}{Sem o 9,93 ($n=11$): média $4{,}48$, mediana $4{,}33$, $s=1{,}71$ e $\mathrm{IQR}=1{,}09$. O desvio padrão cai de $2{,}27$ para $1{,}71$; a mediana se mexe só $0{,}19$ h: é uma medida \textbf{robusta}.}
\end{frame}

\begin{frame}[fragile]{Em Python: outliers e escore $z$}
\begin{pycode}
x = amostra["Horas_Redes"].to_numpy()

q1, q3 = stats.scoreatpercentile(x, [25, 75])
lo, hi = q1 - 1.5*stats.iqr(x), q3 + 1.5*stats.iqr(x)   # 1.69, 7.82
x[(x < lo) | (x > hi)]                  # array([9.93])
stats.zscore(x, ddof=1).max()           # 2.20

stats.tmean(x, limits=(lo, hi))         # 4.48 (sem o outlier)
stats.tstd(x, limits=(lo, hi))          # 1.71
\end{pycode}
\tip{emascsky}{\code{limits=(lo, hi)} faz o SciPy ignorar os valores fora dos limites. A máscara \code{(x < lo) | (x > hi)} é do NumPy.}
\end{frame}

% ============================================================
\setstage{VISUALIZAÇÃO E EXPLORAÇÃO}
\section{Visualização e exploração de dados}

\begin{frame}[fragile]{Escolhendo o gráfico}
\motive{Números resumem; gráficos revelam.}
\centering\small
\begin{tabular}{llll}
\toprule
\textbf{Pergunta} & \textbf{Gráfico} & \textbf{Na nossa base} & \textbf{Função} \\
\midrule
Frequência de cada categoria? & barras & \code{Escolaridade} & \code{plt.bar} \\
Formato da distribuição? & histograma & \code{Horas\_Redes} & \code{plt.hist} \\
Mediana, quartis e outliers? & boxplot & \code{Horas\_Redes} & \code{plt.boxplot} \\
Duas variáveis se relacionam? & dispersão & \code{Horas\_Redes}, \code{Saude\_Mental} & \code{plt.scatter} \\
\bottomrule
\end{tabular}
\vspace{4pt}
\begin{pycode}
fig, ax = plt.subplots()      # cria a figura e os eixos
ax.hist(x)                    # troque hist por boxplot, scatter, bar...
plt.show()
\end{pycode}
\end{frame}

\begin{frame}{Histograma}
\only<1>{\motive{Passo 1. Os doze valores sobre a reta.}}
\only<2>{\motive{Passo 2. Dividir o intervalo em faixas de mesma largura (2 h) e contar os alunos de cada faixa.}}
\only<3>{\motive{Passo 3. Uma barra por faixa, com altura igual à contagem, coladas uma na outra.}}
\only<1>{\figbox{3.9cm}{hist_1}}\only<2>{\figbox{3.9cm}{hist_2}}\only<3>{\figbox{3.9cm}{hist_3}}\par\vspace{2pt}
\only<1>{\tip{emascsky}{Dados: $1{,}92;\ 2{,}04;\ 3{,}61;\ 4{,}12;\ 4{,}17;\ 4{,}33;\ 4{,}71;\ 4{,}91;\ 5{,}00;\ 7{,}10;\ 7{,}39;\ 9{,}93$.}}
\only<2>{\tip{emascsky}{Faixas $[0,2)$, $[2,4)$, $[4,6)$, $[6,8)$ e $[8,10]$. Contagens: $1$; $2$; $6$; $2$; $1$.}}
\only<3>{\tip{emascorange}{As barras são coladas porque a variável é \textbf{contínua}. Faixas estreitas geram ruído; faixas largas escondem o formato.}}
\end{frame}

\begin{frame}{A mesma variável na população}
\motive{Com 16.000 alunos o formato aparece: um sino quase simétrico (assimetria 0,16).}
\figbox{4.0cm}{hist_pop}\par\vspace{2pt}
\tip{emascsky}{Com $n=12$ não dá para afirmar o formato: o histograma da amostra tem aparência irregular. Quanto maior a amostra, mais ela se parece com a população.}
\end{frame}

\begin{frame}[fragile]{Em Python: histograma}
\begin{pycode}
x = amostra["Horas_Redes"].to_numpy()
faixas = [0, 2, 4, 6, 8, 10]

stats.binned_statistic(x, x, "count", bins=faixas).statistic
# array([1., 2., 6., 2., 1.])

plt.hist(x, bins=faixas, edgecolor="black")
plt.xlabel("horas por dia em redes sociais")
plt.show()
\end{pycode}
\tip{emascsky}{\code{stats.binned\_statistic} conta os alunos por faixa (confere com a contagem à mão); \code{plt.hist} só desenha.}
\end{frame}

\begin{frame}[fragile]{Gráfico de barras}
\begin{columns}[c]
\begin{column}{.52\textwidth}\figbox{4.2cm}{barras}\end{column}
\begin{column}{.44\textwidth}
\passo{1}{Listar as categorias no eixo horizontal, \textbf{em ordem}, pois \code{Escolaridade} é ordinal.}
\passo{2}{Contar as observações de cada uma.}
\passo{3}{Uma barra por categoria, com altura igual à contagem, \textbf{separadas}.}
\end{column}
\end{columns}
\begin{pycode}
df["Escolaridade"].value_counts()[ordem]    # 6083, 4890, 5027
df["Escolaridade"].value_counts()[ordem].plot.bar()
\end{pycode}
\end{frame}

\begin{frame}{Boxplot}
\only<1>{\motive{Passo 1. Os doze valores sobre o eixo.}}
\only<2>{\motive{Passo 2. Calcular $Q_1$, mediana e $Q_3$.}}
\only<3>{\motive{Passo 3. Desenhar a caixa de $Q_1$ a $Q_3$, com um traço na mediana.}}
\only<4>{\motive{Passo 4. Estender os bigodes até os valores dentro dos limites e marcar os demais como pontos.}}
\only<1>{\figbox{4.3cm}{box1}}\only<2>{\figbox{4.3cm}{box2}}\only<3>{\figbox{4.3cm}{box3}}\only<4>{\figbox{4.3cm}{box4}}\par\vspace{2pt}
\only<1>{\tip{emascsky}{Mesmos dados do histograma.}}
\only<2>{\tip{emascsky}{$Q_1=3{,}99$, mediana $=4{,}52$, $Q_3=5{,}53$ (Seção 2).}}
\only<3>{\tip{emascsky}{A largura da caixa é o $\mathrm{IQR}=1{,}53$ e cobre os 50\% centrais dos dados.}}
\only<4>{\tip{emascorange}{Limites: $1{,}69$ e $7{,}82$. Bigodes: do menor valor dentro dos limites (1,92) ao maior (7,39). O 9,93 aparece isolado.}}
\end{frame}

\begin{frame}[fragile]{Em Python: boxplot}
\begin{pycode}
plt.boxplot(x, vert=False)
plt.xlabel("horas por dia em redes sociais")
plt.show()
\end{pycode}
\tip{emascsky}{O \code{plt.boxplot} usa os mesmos quartis interpolados e a mesma regra de $1{,}5\times\mathrm{IQR}$ que calculamos à mão.}
\end{frame}

\begin{frame}[fragile]{Gráfico de dispersão}
\motive{Como desenhar: um ponto $(x_i,y_i)$ para cada aluno.}
\begin{columns}[c]
\begin{column}{.46\textwidth}\figbox{4.6cm}{disp1}\end{column}
\begin{column}{.50\textwidth}
\passo{1}{Um aluno vira um ponto: $x$ = horas em redes, $y$ = saúde mental.}
\passo{2}{Observar a nuvem: \textbf{direção}, \textbf{formato}, \textbf{força} e pontos isolados.}
\begin{pycode}
y = amostra["Saude_Mental"].to_numpy()
plt.scatter(x, y)
plt.xlabel("horas em redes")
plt.ylabel("saúde mental")
\end{pycode}
\end{column}
\end{columns}
\end{frame}

\begin{frame}{Correlação entre variáveis}
\motive{Quando uma variável aumenta, a outra tende a aumentar, diminuir ou nada?}
\figbox{2.7cm}{r_grid}\par\vspace{3pt}
\begin{columns}[T]
\begin{column}{.60\textwidth}
\dmath{\mathrm{cov}(x,y)=\frac{1}{n-1}\sum_{i=1}^{n}(x_i-\bar{x})(y_i-\bar{y})}
\dmath{r=\frac{\sum(x_i-\bar{x})(y_i-\bar{y})}{\sqrt{\sum(x_i-\bar{x})^2\,\sum(y_i-\bar{y})^2}},\quad -1\le r\le1}
\end{column}
\begin{column}{.36\textwidth}
\footnotesize
\begin{tabular}{ll}
\toprule
$r$ & \textbf{Interpretação} \\
\midrule
perto de $+1$ & linear positiva forte \\
perto de $0$ & sem relação linear \\
perto de $-1$ & linear negativa forte \\
\bottomrule
\end{tabular}
\end{column}
\end{columns}
\end{frame}

% ============================================================
\setstage{PROBABILIDADE}
\section{Probabilidade e distribuições}

\begin{frame}[fragile]{Conceito de probabilidade}
\motive{Probabilidade quantifica a incerteza com um número entre 0 e 1.}
\begin{columns}[T]
\begin{column}{.31\textwidth}\card{emascsky}{c1}{Experimento aleatório}{Resultado imprevisível.\par\vspace{4pt}{\scriptsize\color{softgray}Sortear um aluno do CSV.}}\end{column}
\begin{column}{.31\textwidth}\card{emascteal}{c1}{Espaço amostral $\Omega$}{Todos os resultados possíveis.\par\vspace{4pt}{\scriptsize\color{softgray}Os 16.000 alunos.}}\end{column}
\begin{column}{.31\textwidth}\card{emascorange}{c1}{Evento}{Subconjunto de interesse.\par\vspace{4pt}{\scriptsize\color{softgray}$A$: mais de 6 h em redes.}}\end{column}
\end{columns}
\vspace{4pt}
\centering
\dmathf{P(A)=\frac{\text{casos favoráveis}}{\text{casos possíveis}}=\frac{4.356}{16.000}\approx0{,}272,\qquad 0\le P(A)\le1}
\begin{pycode}
A = df["Horas_Redes"] > 6          # evento A
A.mean()                           # P(A) = 0.272 (proporção de True)
\end{pycode}
\end{frame}

\begin{frame}{Regras essenciais}
\begin{columns}[c]
\begin{column}{.42\textwidth}\figbox{4.4cm}{venn_regras}\end{column}
\begin{column}{.54\textwidth}
\small
\textbf{Complementar:}\ $P(A^c)=1-P(A)$
\vspace{2pt}

\textbf{União:}\ $P(A\cup B)=P(A)+P(B)-P(A\cap B)$
\vspace{2pt}

\textbf{Condicional:}\ $P(A\mid B)=\dfrac{P(A\cap B)}{P(B)}$
\vspace{2pt}

\textbf{Independência:}\ $P(A\cap B)=P(A)\,P(B)$
\end{column}
\end{columns}
\end{frame}

\begin{frame}{Exemplo: gênero e uso de redes}
\only<1>{\motive{Na amostra de 12: $A$ = Feminino (6), $B$ = mais de 6 h em redes (3). Apenas 1 aluna está nos dois.}}
\only<2>{\motive{União: quem é do gênero feminino ou usa mais de 6 h em redes.}}
\only<3>{\motive{Condicional: entre as alunas, a fração que usa mais de 6 h.}}
\begin{columns}[c]
\begin{column}{.48\textwidth}\only<1>{\figbox{4.4cm}{venn1}}\only<2>{\figbox{4.4cm}{venn2}}\only<3>{\figbox{4.4cm}{venn3}}\end{column}
\begin{column}{.48\textwidth}
\only<1>{\passo{1}{Contar cada região.}\small\vspace{2pt} Só A: $6-1=5$. Os dois: $1$. Só B: $3-1=2$. Nenhum: $12-8=4$.}
\only<2>{\passo{2}{Aplicar a regra da união.}\dmath{P(A\cup B)=\frac{6}{12}+\frac{3}{12}-\frac{1}{12}=\frac{8}{12}\approx0{,}67}}
\only<3>{\passo{3}{Probabilidade condicional.}\dmath{P(B\mid A)=\frac{1/12}{6/12}=\frac{1}{6}\approx0{,}17}\passo{4}{Comparar com $P(B)$.}\dmath{P(B)=\frac{3}{12}=0{,}25\neq0{,}17}\small Na amostra, $A$ e $B$ parecem dependentes.}
\end{column}
\end{columns}
\end{frame}

\begin{frame}[fragile]{Em Python}
\begin{pycode}
A = amostra["Genero"] == "Feminino"
B = amostra["Horas_Redes"] > 6
(A | B).mean()                     # 0.667  P(A união B)
(A & B).sum() / A.sum()            # 0.167  P(B | A)
B.mean()                           # 0.25   P(B)

Ap, Bp = df["Genero"] == "Feminino", df["Horas_Redes"] > 6
(Ap & Bp).sum() / Ap.sum()         # 0.272  P(B | A) na população
Bp.mean()                          # 0.272  P(B) na população
\end{pycode}
\tip{emascorange}{Na população, $P(B\mid A)\approx P(B)$: os eventos são (quase) independentes. Os 12 alunos sugeriam o contrário: é o \textbf{acaso da amostra}.}
\end{frame}

\begin{frame}[fragile]{Distribuição normal}
\motive{Forma de sino, simétrica em torno da média $\mu$ e controlada por $\sigma$.}
\dmath{f(x)=\frac{1}{\sigma\sqrt{2\pi}}\exp\!\left(-\frac{(x-\mu)^2}{2\sigma^2}\right),\qquad X\sim N(\mu,\sigma^2)}
\begin{columns}[c]
\begin{column}{.52\textwidth}\figbox{4.0cm}{normal_pop}\end{column}
\begin{column}{.46\textwidth}
\small A curva usa a média e o desvio padrão da população: $\mu=4{,}54$ e $\sigma=2{,}40$. A altura do sino é $f(\mu)=0{,}166$.
\begin{pycode}
P = df["Horas_Redes"].to_numpy()
mu = stats.tmean(P)
sigma = stats.tstd(P)
N = stats.norm(mu, sigma)
N.pdf(mu)
\end{pycode}
\end{column}
\end{columns}
\end{frame}

\begin{frame}[fragile]{Probabilidade é área sob a curva}
\motive{$P(a\le X\le b)$ é a área entre $a$ e $b$: a diferença entre duas funções de distribuição $F$.}
\dmath{P(a\le X\le b)=F(b)-F(a)}
\begin{columns}[c]
\begin{column}{.52\textwidth}\figbox{3.7cm}{normal_area}\end{column}
\begin{column}{.46\textwidth}
\small Modelo: $0{,}467$. Dados: $0{,}457$. Eles concordam. O modelo falha no extremo esquerdo: 3,6\% dos alunos têm exatamente 0 h.
\begin{pycode}
N.cdf(6) - N.cdf(3)
((P >= 3) & (P <= 6)).mean()
\end{pycode}
\end{column}
\end{columns}
\end{frame}

\begin{frame}{Distribuição binomial}
\motive{Conta o número de sucessos em um número fixo de tentativas independentes.}
\dmath{X\sim\mathrm{Binomial}(n,p),\qquad P(X=k)=\binom{n}{k}p^{k}(1-p)^{n-k}}
\begin{columns}[T]
\begin{column}{.48\textwidth}
\card{emascsky}{b1}{Variáveis}{\small
$n$: número de tentativas (fixo);\par
$p$: probabilidade de sucesso em cada tentativa;\par
$k$: número de sucessos ($k=0,1,\dots,n$).}
\end{column}
\begin{column}{.48\textwidth}
\card{emascorange}{b1}{Que probabilidades calcula}{\small
``Qual a chance de exatamente $k$ sucessos em $n$ tentativas?''\par\vspace{3pt}
\textbf{Ex.:} sorteando $n=5$ alunos do CSV, $p=0{,}272$ de usar mais de 6 h em redes. $P(X=2)=?$}
\end{column}
\end{columns}
\vspace{3pt}
\dmath{E[X]=np,\qquad \operatorname{Var}(X)=np(1-p)}
\end{frame}

\begin{frame}{Distribuição de Poisson}
\motive{Conta eventos raros em um intervalo fixo de tempo, espaço ou volume.}
\dmath{X\sim\mathrm{Poisson}(\lambda),\qquad P(X=k)=\frac{e^{-\lambda}\lambda^{k}}{k!}}
\begin{columns}[T]
\begin{column}{.48\textwidth}
\card{emascteal}{p1}{Variáveis}{\small
$\lambda>0$: taxa média de eventos por intervalo;\par
$k$: número de eventos observados ($k=0,1,2,\dots$).}
\end{column}
\begin{column}{.48\textwidth}
\card{emascorange}{p1}{Que probabilidades calcula}{\small
``Qual a chance de $k$ eventos em um intervalo?''\par\vspace{3pt}
\textbf{Ex.:} só 0,7\% dos alunos usam IA mais de 7 h por dia. Em grupos de 200 alunos, $\lambda=200\cdot0{,}0069\approx1{,}39$. $P(X=3)=?$}
\end{column}
\end{columns}
\vspace{3pt}
\dmath{E[X]=\lambda,\qquad \operatorname{Var}(X)=\lambda}
\end{frame}

\begin{frame}[fragile]{Em Python: binomial e Poisson}
\begin{pycode}
p = (df["Horas_Redes"] > 6).mean()          # 0.272
stats.binom.pmf(2, n=5, p=p)                # 0.286
stats.binom.cdf(2, n=5, p=p)                # 0.872  P(X <= 2)

lam = 200 * (df["Horas_IA"] > 7).mean()     # 1.39
stats.poisson.pmf(3, mu=lam)                # 0.111
stats.poisson.cdf(3, mu=lam)                # 0.948  P(X <= 3)
\end{pycode}
\tip{emascsky}{\code{pmf} dá $P(X=k)$; \code{cdf} dá $P(X\le k)$. As duas distribuições são discretas: $k$ inteiro.}
\end{frame}

% \resrow{cor}{nº}{seção}{pergunta}{ferramentas}{funções Python}{resultado no exemplo}
\newcommand{\resrow}[7]{%
\begin{tcolorbox}[enhanced, colback=emascmist, colframe=#1, boxrule=0pt, arc=3pt,
left=6pt, right=6pt, top=2pt, bottom=2pt, before skip=2.5pt, after skip=0pt,
borderline west={3.5pt}{0pt}{#1}]
\begin{minipage}[c]{.35\linewidth}
{\bfseries\color{#1!85!black}#2.\ #3}\par{\scriptsize\itshape\color{softgray}#4}
\end{minipage}\hfill
\begin{minipage}[c]{.32\linewidth}\raggedright
{\footnotesize #5}\par\vspace{1pt}{\scriptsize #6}
\end{minipage}\hfill
\begin{minipage}[c]{.25\linewidth}\raggedleft
{\footnotesize\bfseries\color{#1!85!black}#7}
\end{minipage}
\end{tcolorbox}}

\begin{frame}{Resumo do minicurso}
\motive{Quatro etapas, uma mesma variável: \code{Horas\_Redes}, da amostra de 12 à população de 16.000.}
\begin{tcolorbox}[enhanced, colback=white, colframe=white, boxrule=0pt, arc=0pt, left=6pt, right=6pt, top=0pt, bottom=0pt,
before skip=0pt, after skip=0pt, borderline west={3.5pt}{0pt}{white}]
{\scriptsize\color{softgray}\bfseries
\begin{minipage}[c]{.35\linewidth}\ \end{minipage}\hfill
\begin{minipage}[c]{.32\linewidth}\raggedright FERRAMENTAS\end{minipage}\hfill
\begin{minipage}[c]{.25\linewidth}\raggedleft NO NOSSO EXEMPLO\end{minipage}}
\end{tcolorbox}
\resrow{emascindigo}{1}{Introdução}{Que dados tenho e em que condições estão?}
{população, amostra, tipos de variável}{\code{read\_csv}, \code{sample}, \code{dtypes}}{$N=16.000$ (população)\par $n=12$ (amostra)}
\resrow{emascsky}{2}{Descritiva}{Onde estão os dados e o quanto variam?}
{$\bar{x}$, mediana, $s$, quartis, $z$}{\code{tmean}, \code{tstd}, \code{zscore}}{$\bar{x}=4{,}94$ h (média)\par $4{,}52$ h (mediana)}
\resrow{emascteal}{3}{Visualização}{Qual o formato? Há valores incomuns?}
{histograma, boxplot, dispersão, $r$}{\code{hist}, \code{boxplot}, \code{scatter}, \code{pearsonr}}{$r=-0{,}50$ (amostra)\par $r=-0{,}40$ (população)}
\resrow{emascamber}{4}{Probabilidade}{Quão provável ou incomum é um valor?}
{$P(A\mid B)$, normal, binomial, Poisson}{\code{norm}, \code{binom}, \code{poisson}}{$P(X>9{,}93)$\par $\approx 1{,}25\%$}
\end{frame}


% -------------------- SLIDE FINAL --------------------
{\darkcanvas
\begin{frame}[plain, noframenumbering]
\begin{pagetikz}
\bigsigma{11.6}{3.0}{3.4}{0.06}
\lammotif{9.2}{2.3}{1.3}{emasccyan}{0.5}
\node[anchor=west, text=white, font=\fontsize{44}{48}\bfseries\selectfont] at (1.25,6.3) {\bigtext Obrigado!};
\fill[emascorange] (1.3,5.35) rectangle (2.9,5.41);
\fill[emasccyan] (3.0,5.35) rectangle (3.6,5.41);
\node[anchor=west, text=emasccyan, font=\fontsize{18}{20}\bfseries\selectfont] at (1.25,4.7) {Dúvidas?};
\node[anchor=north west, text=white, font=\fontsize{11}{14}\bfseries\selectfont, align=left] at (1.25,4.0)
{Gabriel Carvalho\enspace\textcolor{emasccyan}{\textbullet}\enspace João Vitor\enspace\textcolor{emasccyan}{\textbullet}\enspace Kaue Santos \\
Phelipe Melo\enspace\textcolor{emasccyan}{\textbullet}\enspace Rafael Silva};
\node[anchor=north west, text=white!78!emascindigo, font=\fontsize{9}{12}\selectfont, align=left] at (1.25,2.85)
{Liga Acadêmica Matemática \& Empresa -- UFS \\ \ligaemail};
\fill[white] (0,0) rectangle (16,1.65);
\fill[emascorange] (0,1.65) rectangle (5.3,1.69);
\shade[left color=emasccyan, right color=emascsky] (5.3,1.65) rectangle (16,1.69);
\node at (4.6,0.83)  {\includegraphics[height=1.25cm]{figuras/logo_ufs.png}};
\node at (8.0,0.83)  {\includegraphics[height=1.3cm]{figuras/logo_emasc.png}};
\node at (11.4,0.83) {\includegraphics[height=1.25cm]{figuras/logo_liga.png}};
\draw[emascline, line width=0.5pt] (6.3,0.35)--(6.3,1.3);
\draw[emascline, line width=0.5pt] (9.7,0.35)--(9.7,1.3);
\end{pagetikz}
\end{frame}}

\end{document}
